\Needspace{5\baselineskip}
\SourceLeaf{429}%
\begin{center}10.\end{center}
\noindent\hspace*{1.2em}Bevor wir den eben gefundenen Ausdruck weiter umformen, wollen wir einige Bemerkungen dazu machen.

Einer Normalen auf die Kurve in ihrer Ebene entsprechen zwei Richtungen auf der Kugeloberfläche, je nachdem, auf welche Seite der Kurve man sie zieht. Die Richtung, nach der die Kurve \textls[100]{konkav} ist, wird durch $\lambda'$ bezeichnet, die andere durch den gegenüberliegenden Punkt der Kugeloberfläche. Diese beiden Punkte sind ebenso wie $L$ und $\mathfrak{L}$ von $\lambda$ um $90^{0}$ entfernt und liegen daher auf einem Großkreis. Da auch $\mathfrak{L}$ von $\lambda$ um $90^{0}$ entfernt ist, gilt $\mathfrak{L}L=90^{0}-L\lambda'$ oder $=L\lambda'-90^{0}$ und daher
\[
\cos L\lambda'=\pm\sin\mathfrak{L}L,
\]
wobei $\sin\mathfrak{L}L$ notwendig positiv ist. Da wir in unserer Untersuchung $r$ als positiv betrachten, hat $\cos L\lambda'$ dasselbe Vorzeichen wie
\[
Z(\xi\xi T+2\xi\eta U+\eta\eta V)
\]
und ein positiver Wert der letzteren Größe zeigt also an, dass $L\lambda'$ kleiner als $90^{0}$ ist oder dass die Kurve nach der Seite hin konkav ist, auf der die Projektion der Flächennormalen in die Ebene liegt. Ein negativer Wert zeigt dagegen an, dass die Kurve nach dieser Seite hin konvex ist. Wir können daher auch allgemein
\[
\frac1r=\frac{Z(\xi\xi T+2\xi\eta U+\eta\eta V)}{\sin\mathfrak{L}L}
\]
setzen, wenn wir den Krümmungsradius im ersten Fall als positiv und im anderen als negativ betrachten wollen. $\mathfrak{L}L$ ist nun der Winkel, den unsere Schnittebene mit der Tangentialebene der gekrümmten Fläche bildet. Man sieht, dass in verschiedenen Schnittebenen durch denselben Punkt und dieselbe Tangente die Krümmungsradien proportional zum Sinus der Neigung sind. Wegen dieser einfachen Beziehung wollen wir uns fortan auf den Fall beschränken, in dem dieser Winkel ein rechter ist und die Schnittebene daher durch die Normale auf die gekrümmte\space%
